% ==STYLE-BEGIN==
\documentclass[UTF8,11pt]{ctexart}
\usepackage{amsmath,amssymb,amsthm}
\usepackage[a4paper,margin=2.4cm,top=2.3cm,bottom=2.3cm]{geometry}
\usepackage{fancyhdr}
\usepackage{enumitem}
\usepackage{hyperref}
\hypersetup{hidelinks}

% ---- 版式 ----
\linespread{1.15}
\setlength{\parskip}{0.35em}
\setlength{\parindent}{1.6em}
\setlist{nosep,leftmargin=2em}

% ---- 页眉 ----
\pagestyle{fancy}
\fancyhf{}
\fancyhead[L]{\small\coursename}
\fancyhead[R]{\small\thepage}
\renewcommand{\headrulewidth}{0.4pt}

% ---- 定理环境：中文用正体，标题加粗，按小节共用一条流水号 ----
\newtheoremstyle{cnthm}{0.7em}{0.7em}{\normalfont}{}{\bfseries}{}{0.5em}{}
\theoremstyle{cnthm}
\newtheorem{thm}{定理}[section]
\newtheorem{lem}[thm]{引理}
\newtheorem{prop}[thm]{命题}
\newtheorem{cor}[thm]{推论}
\newtheorem{defn}[thm]{定义}
\newtheorem{rem}[thm]{注}
\newtheorem{exa}[thm]{例}
\newtheorem{conj}[thm]{猜想}
\renewcommand{\proofname}{\bfseries 证明}
\providecommand{\coursename}{}
% ==STYLE-END==
\renewcommand{\coursename}{高等代数\,·\,1104}
\xeCJKDeclareCharClass{CJK}{"2150 -> "218F} % 让 Ⅰ Ⅱ Ⅲ 走中文字体
\title{高等代数 1104}
\date{}
\begin{document}
\maketitle

\section{$m$ 维向量空间(回顾)}

\begin{itemize}
\item $K$ 是数域, $K^m=\{(a_1,a_2,\cdots,a_m)^{\mathrm T}\mid a_i\in K,\ i=1,\cdots,m\}$; $(K^m,+,\cdot)$ 是线性空间 (满足八条性质);
\item \textbf{线性组合}: $\sum_{i=1}^{s}k_i\alpha_i$ 称为 $\alpha_1,\cdots,\alpha_s$ 的线性组合;
\item \textbf{线性表示}: 若 $\beta=\sum_{i=1}^{s}k_i\alpha_i$, 称 $\beta$ 可被 $\alpha_1,\cdots,\alpha_s$ 线性表示;
\item \textbf{线性相关/无关}: 看方程 $x_1\alpha_1+\cdots+x_s\alpha_s=0$ 是否有非零解. 它是齐次线性方程组, 其系数矩阵为
\[
A=(\alpha_1,\cdots,\alpha_s).
\]
\end{itemize}

\begin{exa}
$K^m$ 中任意 $(m+1)$ 个向量 $\alpha_1,\cdots,\alpha_{m+1}$ 一定线性相关.

理由: $A=(\alpha_1,\cdots,\alpha_{m+1})$ 的行阶梯形矩阵一定有些列没有主元 (列数 $m+1$ 大于行数 $m$), 故必有非零解. 反过来, $K^m$ 中可以找到 $m$ 个线性无关的向量.
\end{exa}

\begin{prop}
$\alpha_1,\cdots,\alpha_s$ 线性相关 $\iff$ 存在某个 $\alpha_j$, s.t. $\alpha_j$ 能被其余向量线性表示 (即向量组中有``多余''的向量).
\end{prop}

\section{向量组的秩}

\begin{defn}[线性等价]
$K^m$ 中有向量组 $\alpha_1,\cdots,\alpha_s$ (记为 (Ⅰ)) 及 $\beta_1,\cdots,\beta_t$ (记为 (Ⅱ)). 如果 (Ⅰ) 中每个向量能被 (Ⅱ) 线性表示, 且 (Ⅱ) 中每个向量能被 (Ⅰ) 线性表示, 则称 (Ⅰ) 和 (Ⅱ) \textbf{线性等价}, 记作 (Ⅰ)$\sim$(Ⅱ).
\end{defn}

\begin{exa}
$K^3$ 中, $\alpha_1=(1,0,0)^{\mathrm T}$, $\alpha_2=(0,1,0)^{\mathrm T}$ 构成 (Ⅰ); $\beta_1=(1,1,0)^{\mathrm T}$, $\beta_2=(0,2,0)^{\mathrm T}$, $\beta_3=(1,-2,0)^{\mathrm T}$ 构成 (Ⅱ). 则 (Ⅰ) 与 (Ⅱ) 线性等价.
\end{exa}

\begin{thm}[线性等价的性质]
(1) 反身性: (Ⅰ)$\sim$(Ⅰ);

(2) 对称性: (Ⅰ)$\sim$(Ⅱ) $\Rightarrow$ (Ⅱ)$\sim$(Ⅰ);

(3) 传递性: (Ⅰ)$\sim$(Ⅱ), (Ⅱ)$\sim$(Ⅲ) $\Rightarrow$ (Ⅰ)$\sim$(Ⅲ).
\end{thm}

线性等价是``等价关系''.

\begin{lem}
设 $\alpha_1,\cdots,\alpha_s$ (Ⅰ) 及 $\beta_1,\cdots,\beta_t$ (Ⅱ), 且 (Ⅱ) 能被 (Ⅰ) 线性表示. 若 $\gamma$ 能被 (Ⅱ) 线性表示, 则 $\gamma$ 能被 (Ⅰ) 线性表示.
\end{lem}

\begin{proof}
由于 (Ⅱ) 能被 (Ⅰ) 线性表示, 则
\[
\begin{aligned}
\beta_1&=k_{11}\alpha_1+k_{21}\alpha_2+\cdots+k_{s1}\alpha_s,\\
\beta_2&=k_{12}\alpha_1+k_{22}\alpha_2+\cdots+k_{s2}\alpha_s,\\
&\ \ \cdots\cdots\\
\beta_t&=k_{1t}\alpha_1+k_{2t}\alpha_2+\cdots+k_{st}\alpha_s,
\end{aligned}
\]
用矩阵乘法可以写成
\[
(\beta_1,\cdots,\beta_t)=(\alpha_1,\cdots,\alpha_s)
\begin{pmatrix}
k_{11} & k_{12} & \cdots & k_{1t}\\
k_{21} & k_{22} & \cdots & k_{2t}\\
\vdots & \vdots & & \vdots\\
k_{s1} & k_{s2} & \cdots & k_{st}
\end{pmatrix}
=(\alpha_1,\cdots,\alpha_s)K .
\]
又 $\gamma=a_1\beta_1+\cdots+a_t\beta_t$, 用矩阵乘法可以写成
\[
\gamma=(\beta_1,\cdots,\beta_t)\begin{pmatrix}a_1\\a_2\\\vdots\\a_t\end{pmatrix},
\]
于是
\[
\gamma=(\alpha_1,\cdots,\alpha_s)\,K\begin{pmatrix}a_1\\a_2\\\vdots\\a_t\end{pmatrix},
\]
即有
\[
\gamma=\Bigl(\sum_{j=1}^{t}k_{1j}a_j\Bigr)\alpha_1+\cdots+\Bigl(\sum_{j=1}^{t}k_{sj}a_j\Bigr)\alpha_s
=\sum_{i=1}^{s}\Bigl(\sum_{j=1}^{t}k_{ij}a_j\Bigr)\alpha_i ,
\]
所以 $\gamma$ 能被 $\alpha_1,\cdots,\alpha_s$ 线性表示.
\end{proof}

\section{极大线性无关部分组}

\begin{defn}
设 $\alpha_{i_1},\cdots,\alpha_{i_r}$ (Ⅱ) 是 $\alpha_1,\cdots,\alpha_s$ (Ⅰ) 的一个部分组. 如果

(1) (Ⅰ) 和 (Ⅱ) 线性等价;

(2) (Ⅱ) 线性无关,

则称 (Ⅱ) 是 (Ⅰ) 的\textbf{极大线性无关部分组} (简称\textbf{极大无关组}).
\end{defn}

\begin{exa}
$\beta_1=(1,1,0)^{\mathrm T}$, $\beta_2=(0,2,0)^{\mathrm T}$, $\beta_3=(1,-2,0)^{\mathrm T}$. 因为
\[
\beta_3=\beta_1-\frac32\beta_2,
\]
所以 $\beta_1,\beta_2$ 是极大无关组; 同理 $\beta_1,\beta_3$ 也是, 且
\[
\beta_2=\frac23\beta_1-\frac23\beta_3 .
\]
可见\textbf{极大无关组不唯一, 但向量个数相同} (本例都是 $2$).
\end{exa}

\begin{thm}
设 $\alpha_1,\cdots,\alpha_s$ (Ⅰ), $\beta_1,\cdots,\beta_t$ (Ⅱ). 如果 (Ⅰ) 能被 (Ⅱ) 线性表示, 且 $s>t$, 则 (Ⅰ) 一定线性相关.
\end{thm}

\begin{proof}
由于 (Ⅰ) 能被 (Ⅱ) 线性表示, 则
\[
(\alpha_1,\cdots,\alpha_s)=(\beta_1,\cdots,\beta_t)
\begin{pmatrix}
k_{11} & k_{12} & \cdots & k_{1s}\\
k_{21} & k_{22} & \cdots & k_{2s}\\
\vdots & \vdots & & \vdots\\
k_{t1} & k_{t2} & \cdots & k_{ts}
\end{pmatrix}
=(\beta_1,\cdots,\beta_t)K .
\]
考虑方程
\[
x_1\alpha_1+\cdots+x_s\alpha_s=0,
\]
其等价于
\[
(\alpha_1,\cdots,\alpha_s)\begin{pmatrix}x_1\\\vdots\\x_s\end{pmatrix}=0
\iff
(\beta_1,\cdots,\beta_t)\,K\begin{pmatrix}x_1\\\vdots\\x_s\end{pmatrix}=0 .
\]
即得: 只要 $x\neq0$ 满足 $Kx=0$, 就有 $x_1\alpha_1+\cdots+x_s\alpha_s=0$. 现考虑方程
\[
K\begin{pmatrix}x_1\\\vdots\\x_s\end{pmatrix}=0 ,
\]
这是 $t$ 个方程、$s$ 个未知量的齐次线性方程组. 由于 $s>t$, 它必有非零解, 故存在不全为零的 $x_1,\cdots,x_s$ 使 $x_1\alpha_1+\cdots+x_s\alpha_s=0$, 即 $\alpha_1,\cdots,\alpha_s$ 线性相关.
\end{proof}

\begin{thm}
设 $\alpha_1,\cdots,\alpha_s$ (Ⅰ), 且 $\alpha_{i_1},\cdots,\alpha_{i_t}$ (Ⅱ) 与 $\beta_{j_1},\cdots,\beta_{j_r}$ (Ⅲ) 都是 (Ⅰ) 的极大无关组, 则 $t=r$.
\end{thm}

\begin{cor}
设 $\alpha_1,\cdots,\alpha_s$ (Ⅰ) 与 $\beta_1,\cdots,\beta_t$ (Ⅱ). 如果 (Ⅰ) 能被 (Ⅱ) 线性表示, 且 (Ⅰ) 线性无关, 则 $s\leq t$.
\end{cor}

\begin{defn}[向量组的秩]
\[
\operatorname{rank}(\alpha_1,\cdots,\alpha_s)\triangleq \alpha_1,\cdots,\alpha_s \text{ 任一极大无关组的向量个数}.
\]
\end{defn}

\section{筛选法选出极大无关组}

设 $\alpha_1,\cdots,\alpha_s$. 不妨设 $\alpha_1\neq0$;

若 $\alpha_2$ 能被 $\alpha_1$ 线性表示, 则去掉 $\alpha_2$, 否则保留 $\alpha_2$;

重复上述操作: 对于 $\alpha_j$, 如果 $\alpha_j$ 能被前面\textbf{保留}的向量线性表示, 则去掉 $\alpha_j$, 否则保留 $\alpha_j$, 直到选择完成. 此时得到向量组 $\alpha_{i_1},\cdots,\alpha_{i_r}$.

\begin{proof}[为什么它是极大无关组]
被去掉的每个向量都能被保留下来的向量线性表示, 故两者线性等价. 若 $\alpha_{i_1},\cdots,\alpha_{i_r}$ 线性相关, 则存在不全为零的 $k_1,\cdots,k_r$ 使
\[
k_1\alpha_{i_1}+\cdots+k_r\alpha_{i_r}=0 ;
\]
从左往右取第一个不为零的系数 $k_j$, 则
\[
k_1\alpha_{i_1}+\cdots+k_j\alpha_{i_j}=0
\quad\Longrightarrow\quad
\alpha_{i_j}=-\frac{1}{k_j}\bigl(k_1\alpha_{i_1}+\cdots+k_{j-1}\alpha_{i_{j-1}}\bigr),
\]
即 $\alpha_{i_j}$ 能被前面保留的向量线性表示, 与``保留''的规则矛盾. 所以保留的向量组线性无关, 从而是极大无关组.
\end{proof}

\end{document}
